\section{Introduction}
\label{sec:introduction}

A portfolio left untouched does not stay where it was put: differential asset
returns drift realised weights away from the risk--return position that was
chosen~\cite{markowitz1952,sharpe1966}. Rebalancing restores that position, but
every restoration pays spread, brokerage and impact costs and liquidates
positions whose recent strength may persist~\cite{jegadeesh1993}, so the
decision is a cost--benefit test on the current drift, market state and
transaction cost jointly. Incumbent practice answers with rules that use almost
none of this information: calendar rebalancing acts on a fixed interval and
drift-band rebalancing on a fixed threshold, both blind to the volatility,
drawdown and trend that determine whether a round trip is worth its price. The
applied literature argues this conceptually rather than measuring
it~\cite{adedoyin_ai,kumari2024}; measured against forward, net-of-cost outcomes
on our data, a 90-day calendar rule and a drift-band rule attain test AUC
$0.5167$ and $0.5146$, both close to chance. Fig.~\ref{fig:f1} shows the
system in which we replace such a rule with a learned detector.

The data this decision needs cannot reasonably be centralised: holdings
disclose wealth, leverage and trading history and remain identifying long
after any label is stripped, while a single participant produces too few
rebalancing states to learn a market-wide regularity alone. Federated
averaging resolves this cleanly for a linear detector, since its
parameters --- a twelve-element weight vector and intercept --- are exactly
the object federated averaging is defined over~\cite{mcmahan2017}, needing no
distillation step. Transmitting updates instead of holdings is data
minimisation, not a privacy guarantee; no differential privacy or secure
aggregation is implemented anywhere in this paper.

Supervised rebalancing has no naturally observable label, and the expedient
response is to label historical states with the output of the very heuristic
the model should replace. We call the resulting pathology \emph{rule-oracle
degeneracy}: the label is a deterministic function of the features the rule
reads, so the classifier decompiles the rule instead of learning an economic
decision. Section~\ref{sec:implementation} shows this is large enough to be
mistaken for success: on one fixed corpus, holding everything but the label
function constant, the same pipeline reaches test AUC $0.9890$ against
heuristic labels and $0.6489$ against a forward, net-of-cost backtest.

\subsection{Contributions}
\begin{itemize}
\item \textbf{C1: Backtest-derived labels for rebalancing-necessity detection.}
Each label is the sign of a forward 21-trading-day, net-of-cost
certainty-equivalent gain (Section~\ref{sec:methodology}), evaluated on
$3{,}706$ NIFTY~50 sessions (2011-08-01--2026-07-31) at $0.35\%$ round-trip
cost, over $30{,}000$ generated states with structural positive rate $0.5009$.
\item \textbf{C2: A quantified measurement of rule-oracle degeneracy.} A
matched control varying only the label function measures a gap of $0.3401$
test AUC attributable to the label, not the learner (Section~\ref{sec:implementation}).
\item \textbf{C3: A hardened federated aggregation protocol.} Updates are
$L_2$-clipped and combined by a sample-weighted mean, gated by a minimum
participant count, an upload cap and a per-account throttle, with atomic
activation. It is implemented and unit-tested but disabled by default and not
experimentally evaluated here.
\item \textbf{C4: A cross-runtime parity contract} that pins feature order and
scaler constants by version, so pooling is well defined and mismatched
contracts are rejected rather than averaged, to $10^{-12}$/$10^{-5}$ agreement.
\item \textbf{C5: An explicit model-selection trade.} The linear detector
surrenders $0.0301$ test AUC to a random forest and is chosen anyway because
its parameters are averageable and its probabilities are calibrated (ECE
$0.0214$); the deficit is reported, not suppressed.
\end{itemize}

The evidence is correspondingly narrow: the cross-section is generated rather
than observed, no path-based backtest or federated experiment was run, and
none of this constitutes investment advice. Section~\ref{sec:literature}
surveys the related literature; Section~\ref{sec:methodology} presents
detection and federation as equally weighted halves of one design;
Section~\ref{sec:implementation} reports the results; Section~\ref{sec:conclusion}
concludes.
